\section{Models}
\label{models}

Nature (Real world, Universe) is unique and exists independently of us. The observer who appears in physical
literature is nothing more than figure of speech. `Moon exists even nobody looks at it' and so does electron.
Real word objects behave according to Laws of Nature. Contrary to human laws the Laws of Nature could never be violated.
When we read something like 'this device violates laws of physics' it means usually misunderstanding or clickbait (probably both). Discovering of these laws is the goal of physics. 

To achieve this goal
we build Ideal world and populate it by models. These models emerge as result of  observations (experiments, measurements, experience), 
their elements appear as ideal images of real objects. Using instruments, 
which may be as simple as measuring tape or lever scaling or as complex as electronic microscope or collider,
we get facts upon which model is based and updated.
It should be mentioned that our instruments are part of Real world  and hence described by models\footnote{so a person trying 
to disprove say quantum mechanics would have serious problem since every modern instrument is based actually on quantum 
mechanics.}

    When constructing a model, the following idealization is
   made: certain facts which are only known with a certain degree of
   probability or with a certain degree of accuracy, are considered to be
   "absolutely" correct and are accepted as "axioms". The sense of this
   "absoluteness" lies precisely in the fact that we allow ourselves to
   use these "facts" according to the rules of formal logic, in the
   process declaring as "theorems" all that we can derive from them \cite{Arnold-mt}.

   Events which are  highly probable in Real world become certain in ideal one, plausible deductions become proofs  and so on. Even the statement that results of measurement are real numbers is actually idealization: no instrument is able to produce output with infinite number of digits, hence experimental results are always rational numbers.

Model should be {\em correct}. It should not contradict already known facts. For example no
mechanical model, where energy and moment do not conserve in elastic collisions, can
be correct.

Model should be {\em consistent}. It should  not contradict itself. No statement can be proven to
be both false and true, though its status can be unknown or even undefined.

Model should be {\em efficient}. It should allow to obtain  specific results,
at least in some particular case (e.g. using perturbation theory). 
Let us call such model {\em reasonable}.

 When reasonable model is constructed, 
it is investigated and conclusions are mapped back to Real world as predictions which 
should be tested by  experiment. 
If we get  accurate description of reality the model remains {\em reasonable} and can be foundation 
for other models, in particular those describing our instruments, if not it  should be replaced or improved in some way or another.
If the model fails to produce results to be experimentally checked, it simply makes no sense.

Let's assume that we have reasonable (mathematical) model. How can
we obtain correct conclusions from it. Here mathematician
and physicist take different paths. Mathematician rigorously
and carefully investigates the model as it is formulated.
Physicist may change or abandon the model if it yields 
 incorrect results (contradicting to known experimental facts).
Nature is too complex to be described by one model.
Physicist often uses more than one model simultaneously and change them on the fly 
without warning or use some conditions which were not stated explicitly. 
These tricks irritate mathematicians and confuse engineers. However they are natural parts of 
physicist's way of thinking. Contrary to what is said in many (good) books the relations 
among models are not hierarchical. Generally it is wrong to say that one physical model is a particular
case of another valid only in restricted domain. 

 These relations resemble cartography \footnote{Or definition of smooth manifold}: 
Earth can not have a plane representation consisting of a single map and 
therefore one needs atlas for covering the whole Earth surface. 
Atlas is a set of maps which together cover the Earth. 
When two maps overlap the  conversion between them should be defined.

So Ideal world can be considered as set of models each of them explain particular aspect
of the Real world. Together they (hopefully) cover all the Real world,  and  (not necessary one-to-one)
 transformation between overlapping models exists.
For example only system of material points interacting via electromagnetic or (and) gravitational
field constitutes the common subject of classical, relativistic and quantum mechanics.

Evidently such things as material point, absolutely rigid body etc.  make sense only in the Ideal world
 and should be mapped with care to Real world. 
On the other hand
some notions like causality,  likelihood,
`experimental support' make sense only and in Real world and need 
special care when mapped to Ideal world. 

Why do we need an idealization? The answer is: to be able to use mathematical machinery, in 
particular (long) chains of reasoning. Imagine using  approximate equality instead exact one.
Then we can easily 'prove' that $ 1 \approx 2 $ with accuracy better than 10 percent
using the chain of equations: $ 1 \approx  1.1$, $1.1 \approx 1.2$,\ldots  $  1.9 \approx 2 $.
Another example: the sequence of 100 almost certain events, having probability $0.99$ will
may occur with probability $0.99^{100} = 0.36$. 

However, in the ways of using mathematical tools, in the very character
of reasoning, theoretical physics possesses certain specific features
that not only prevent it from being regarded simply as a branch of
mathematics, but sometimes even make its methods in some sense directly
opposite to mathematical ones.
 “The mathematician proves, while the physicist convinces.” \cite{Medvedev}

A lot of paradoxes appear when we mix Real and Ideal world. 
Let us take for example  Hempel's raven paradox: ``observing
of red cow increases the likelihood 
that all ravens are black''. 

Indeed suppose that the statement `All ravens are black' is true.
In the form of an 
implication, this can be expressed as: If something is a raven, then it 
is black. Via contraposition, this statement is equivalent to: If something is 
not black, then it is not a raven. So an observation of, say, red cow supports 
the statement that all ravens are black. This looks very strange but actually is quite
straightforward. In formal logic (Ideal world) there is no such thing as experimental support or
likelihood. Statement may be invalidated by the observation of non black raven, 
but can not be confirmed in any way. So contraposition means simply that the 
statement `raven is found which is not black' is equivalent to 
`something not black is found which is raven'. 

On the other hand common sense (Real world) does not allow statements about 
{\em all} ravens. We should say majority or practically all 
and  contraposition  disappears.

Another good example is Schroedinger's cat. In quantum mechanics the (pure) state of the system is described 
by wave function. Nobody knows exact wave function for an object more complicated than hydrogen atom. 
Surely we can use approximate one, but even in that case preparing of the object in pure state
is rather involved task.  Even if somebody can figure out wave function of cat, how to prepare it
in pure state? May be cool it down to $0^\circ K$? Correlations replace causality in Ideal World.
 Life and death exist in Real World only and have no analogs in Ideal World.
